\section{Vectors}

The Cartesian coordinate system allows us to describe points by numbers.  A
second geometric object is needed to describe displacement, direction, and
change.  This object is a vector.

\subsection{Directed Segments and Free Vectors}

Let $A$ and $B$ be points.  The directed segment from $A$ to $B$ is denoted by
\[
\overrightarrow{AB}.
\]
It has a direction, an orientation, and a length.  Two directed segments
represent the same vector when one can be translated onto the other without
rotation or change of length.

\begin{definitionbox}
A \emph{free vector} is an equivalence class of directed segments having the
same length, direction, and orientation.  A particular directed segment is a
representative of the vector.
\end{definitionbox}

Thus a vector is not attached permanently to one point.  It may be translated
parallel to itself anywhere in the plane or in space.

\begin{center}
\begin{tikzpicture}[scale=0.9]
    \draw[axis] (-0.5,0) -- (6.2,0) node[right] {$x$};
    \draw[axis] (0,-0.5) -- (0,4.2) node[above] {$y$};
    \coordinate (A) at (0.8,0.8);
    \coordinate (B) at (3.0,2.4);
    \coordinate (C) at (3.4,1.2);
    \coordinate (D) at (5.6,2.8);
    \draw[subset,->] (A) -- (B) node[midway,above left] {$\mathbf v$};
    \draw[subset,->] (C) -- (D) node[midway,below right] {$\mathbf v$};
    \fill[geometricSubsetColor] (A) circle (2pt) node[below left] {$A$};
    \fill[geometricSubsetColor] (B) circle (2pt) node[above] {$B$};
    \fill[geometricSubsetColor] (C) circle (2pt) node[below] {$C$};
    \fill[geometricSubsetColor] (D) circle (2pt) node[above right] {$D$};
\end{tikzpicture}
\end{center}

The zero vector, denoted by $\mathbf 0$, has length zero.  Its initial and
terminal points coincide.  For every vector $\mathbf v$, the opposite vector
$-\mathbf v$ has the same length and direction but the reverse orientation.

\subsection{Addition and the Parallelogram Rule}

Suppose that representatives of $\mathbf u$ and $\mathbf v$ begin at the same
point.  Translate each vector to the endpoint of the other.  The four directed
segments form a parallelogram.

\begin{definitionbox}
The sum $\mathbf u+\mathbf v$ is the vector represented by the diagonal of this
parallelogram beginning at the common initial point.
\end{definitionbox}

Equivalently, place the initial point of $\mathbf v$ at the terminal point of
$\mathbf u$.  The vector from the initial point of $\mathbf u$ to the terminal
point of the translated $\mathbf v$ is $\mathbf u+\mathbf v$.  This is the
triangle rule.

\begin{center}
\begin{tikzpicture}[scale=0.95]
    \coordinate (O) at (0,0);
    \coordinate (U) at (3.2,0.8);
    \coordinate (V) at (1.1,2.5);
    \coordinate (S) at (4.3,3.3);
    \draw[subset,->] (O) -- (U) node[midway,below] {$\mathbf u$};
    \draw[subset,->] (O) -- (V) node[midway,left] {$\mathbf v$};
    \draw[helper] (U) -- (S);
    \draw[helper] (V) -- (S);
    \draw[very thick,->] (O) -- (S) node[midway,above left] {$\mathbf u+\mathbf v$};
    \fill[geometricSubsetColor] (O) circle (2pt);
\end{tikzpicture}
\end{center}

Vector addition satisfies
\[
\mathbf u+\mathbf v=\mathbf v+\mathbf u,
\qquad
(\mathbf u+\mathbf v)+\mathbf w=\mathbf u+(\mathbf v+\mathbf w),
\]
\[
\mathbf v+\mathbf 0=\mathbf v,
\qquad
\mathbf v+(-\mathbf v)=\mathbf 0.
\]
Subtraction is defined by
\[
\mathbf u-\mathbf v=\mathbf u+(-\mathbf v).
\]

\subsection{Multiplication by a Scalar}

Let $\lambda\in\mathbb R$.

\begin{definitionbox}
The scalar multiple $\lambda\mathbf v$ is the vector whose length is
$|\lambda|$ times the length of $\mathbf v$.  It has the same orientation as
$\mathbf v$ when $\lambda>0$, the opposite orientation when $\lambda<0$, and
is the zero vector when $\lambda=0$.
\end{definitionbox}

Scalar multiplication satisfies
\[
\lambda(\mathbf u+\mathbf v)=\lambda\mathbf u+\lambda\mathbf v,
\qquad
(\lambda+\mu)\mathbf v=\lambda\mathbf v+\mu\mathbf v,
\]
\[
(\lambda\mu)\mathbf v=\lambda(\mu\mathbf v),
\qquad
1\mathbf v=\mathbf v.
\]

Two nonzero vectors are parallel precisely when one is a nonzero scalar
multiple of the other:
\[
\mathbf u\parallel\mathbf v
\iff
\mathbf u=\lambda\mathbf v
\quad\text{for some }\lambda\neq0.
\]

\subsection{Coordinates of a Vector}

Let
\[
A=(x_1,y_1),
\qquad
B=(x_2,y_2).
\]
The vector from $A$ to $B$ has coordinates
\[
\overrightarrow{AB}=(x_2-x_1,\,y_2-y_1).
\]
In three-dimensional space, if
\[
A=(x_1,y_1,z_1),
\qquad
B=(x_2,y_2,z_2),
\]
then
\[
\overrightarrow{AB}=(x_2-x_1,\,y_2-y_1,\,z_2-z_1).
\]

For vectors written in coordinates,
\[
\mathbf u=(u_1,u_2),
\qquad
\mathbf v=(v_1,v_2),
\]
we have
\[
\mathbf u+\mathbf v=(u_1+v_1,\,u_2+v_2),
\qquad
\lambda\mathbf u=(\lambda u_1,\,\lambda u_2).
\]
The analogous formulas hold componentwise in $\mathbb R^3$.

The standard basis vectors are
\[
\mathbf e_1=(1,0),
\qquad
\mathbf e_2=(0,1)
\]
in the plane, and
\[
\mathbf e_1=(1,0,0),
\qquad
\mathbf e_2=(0,1,0),
\qquad
\mathbf e_3=(0,0,1)
\]
in space.  Therefore
\[
(u_1,u_2)=u_1\mathbf e_1+u_2\mathbf e_2
\]
and
\[
(u_1,u_2,u_3)=u_1\mathbf e_1+u_2\mathbf e_2+u_3\mathbf e_3.
\]

\subsection{Linear Combinations, Span, and Bases}

An expression of the form
\[
\lambda_1\mathbf v_1+\cdots+\lambda_k\mathbf v_k
\]
is called a linear combination of the vectors $\mathbf v_1,\ldots,\mathbf v_k$.
The set of all such combinations is their span.

Two vectors in the plane are linearly independent when neither is a scalar
multiple of the other.  In that case they form a basis of the plane: every
vector can be written uniquely as their linear combination.  Similarly, three
vectors in $\mathbb R^3$ form a basis when none of them can be expressed as a
linear combination of the other two.

This language will later describe lines and planes:
\[
P=P_0+t\mathbf v
\]
describes a line, while
\[
P=P_0+s\mathbf u+t\mathbf v
\]
describes a plane when $\mathbf u$ and $\mathbf v$ are linearly independent.

\subsection{Length and Unit Vectors}

The length, or Euclidean norm, of
\[
\mathbf v=(v_1,v_2)
\]
is
\[
\|\mathbf v\|=\sqrt{v_1^2+v_2^2}.
\]
For
\[
\mathbf v=(v_1,v_2,v_3)
\]
in space,
\[
\|\mathbf v\|=\sqrt{v_1^2+v_2^2+v_3^2}.
\]
These formulas follow from the Pythagorean theorem.

A vector of length one is called a unit vector.  Every nonzero vector
$\mathbf v$ determines the unit vector
\[
\widehat{\mathbf v}=\frac{\mathbf v}{\|\mathbf v\|}
\]
in the same direction.

The distance between points is the length of their displacement vector:
\[
d(A,B)=\|\overrightarrow{AB}\|.
\]

\subsection{The Dot Product and Angles}

For vectors in $\mathbb R^2$ or $\mathbb R^3$, define the dot product by
\[
\mathbf u\cdot\mathbf v
=
\sum_i u_i v_i.
\]
Thus, in three dimensions,
\[
(u_1,u_2,u_3)\cdot(v_1,v_2,v_3)
=
u_1v_1+u_2v_2+u_3v_3.
\]

\begin{theorembox}
If $\theta$ is the angle between two nonzero vectors $\mathbf u$ and
$\mathbf v$, then
\[
\mathbf u\cdot\mathbf v
=
\|\mathbf u\|\,\|\mathbf v\|\cos\theta.
\]
Consequently,
\[
\cos\theta
=
\frac{\mathbf u\cdot\mathbf v}
{\|\mathbf u\|\,\|\mathbf v\|}.
\]
\end{theorembox}

The dot product is symmetric and distributive:
\[
\mathbf u\cdot\mathbf v=\mathbf v\cdot\mathbf u,
\qquad
\mathbf u\cdot(\mathbf v+\mathbf w)
=
\mathbf u\cdot\mathbf v+\mathbf u\cdot\mathbf w.
\]
Moreover,
\[
\mathbf v\cdot\mathbf v=\|\mathbf v\|^2.
\]

\subsection{Perpendicularity and Orthogonal Projection}

\begin{definitionbox}
Two nonzero vectors are perpendicular, or orthogonal, when the angle between
them is a right angle.
\end{definitionbox}

Because $\cos\theta=0$ for a right angle,
\[
\mathbf u\perp\mathbf v
\iff
\mathbf u\cdot\mathbf v=0.
\]
This criterion will be used to define normal vectors to lines and planes.

The orthogonal projection of $\mathbf u$ onto a nonzero vector $\mathbf v$ is
\[
\operatorname{proj}_{\mathbf v}\mathbf u
=
\frac{\mathbf u\cdot\mathbf v}{\mathbf v\cdot\mathbf v}\,\mathbf v.
\]
The difference
\[
\mathbf u-\operatorname{proj}_{\mathbf v}\mathbf u
\]
is perpendicular to $\mathbf v$.

\begin{center}
\begin{tikzpicture}[scale=0.95]
    \coordinate (O) at (0,0);
    \coordinate (V) at (5,0);
    \coordinate (U) at (3.2,2.7);
    \coordinate (P) at (3.2,0);
    \draw[subset,->] (O) -- (V) node[right] {$\mathbf v$};
    \draw[subset,->] (O) -- (U) node[above] {$\mathbf u$};
    \draw[very thick,->] (O) -- (P) node[midway,below] {$\operatorname{proj}_{\mathbf v}\mathbf u$};
    \draw[helper] (U) -- (P);
    \draw (P) ++(-0.28,0) -- ++(0,0.28) -- ++(0.28,0);
\end{tikzpicture}
\end{center}

\subsection{The Cross Product in Three Dimensions}

For
\[
\mathbf u=(u_1,u_2,u_3),
\qquad
\mathbf v=(v_1,v_2,v_3),
\]
define
\[
\mathbf u\times\mathbf v
=
(u_2v_3-u_3v_2,\,
 u_3v_1-u_1v_3,\,
 u_1v_2-u_2v_1).
\]

\begin{theorembox}
The vector $\mathbf u\times\mathbf v$ is perpendicular to both $\mathbf u$ and
$\mathbf v$.  Its length is
\[
\|\mathbf u\times\mathbf v\|
=
\|\mathbf u\|\,\|\mathbf v\|\sin\theta,
\]
where $\theta$ is the angle between $\mathbf u$ and $\mathbf v$.
\end{theorembox}

Therefore $\|\mathbf u\times\mathbf v\|$ is the area of the parallelogram
spanned by $\mathbf u$ and $\mathbf v$.  The cross product is zero exactly when
the vectors are parallel.  Its orientation is determined by the right-hand
rule.

The scalar triple product
\[
\mathbf u\cdot(\mathbf v\times\mathbf w)
\]
is the signed volume of the parallelepiped spanned by the three vectors.  It is
zero exactly when the vectors are coplanar.

\begin{summarybox}
Vectors describe displacement and direction.  Their fundamental operations are
addition, scalar multiplication, the dot product, and---in three dimensions---
the cross product.  The parallelogram rule gives vector addition, the dot
product detects angles and perpendicularity, and the cross product produces a
normal vector to a plane.  These constructions provide the language needed for
lines, planes, conic sections, and quadric surfaces.
\end{summarybox}
